\chapter{Modèles d'inventaire}\label{ch:hf:inventory-models}

Un teneur de marché qui vient d'acheter trois fois de suite détient un stock que personne n'a encore demandé. Baisser un peu les deux cotations rend la
prochaine transaction plus susceptible d'être une vente~; baisser trop les cotations revient à céder la fourchette. Avellaneda et Stoikov ont transformé ce compromis en un problème de contrôle dont
la solution approchée tient sur une seule ligne, et l'ont simulé~: à leur aversion au risque de référence, les cotations tenant compte de l'inventaire ont rapporté 5~\% de moins que des cotations symétriques
de même largeur moyenne, avec la moitié de l'écart type. Ce chapitre résout le problème, reproduit leur expérience, traite la variante
avec une \omterm{def:hf:inventory-models:penalty}{pénalité d'inventaire} en continu, la résout exactement, puis en tire la réponse pour un marché à grand tick, où une cotation ne peut se trouver qu'au meilleur niveau ou être absente~: là,
sur les sessions simulées de ce chapitre, le contrôle de l'inventaire transforme un teneur de marché déficitaire en un teneur à l'équilibre.

\section{Le problème de contrôle du teneur de marché}

Ho et Stoll posent le problème en 1981~: un dealer dont les prix dépendent de son inventaire, parce que détenir un inventaire est risqué. La forme moderne conserve les
ingrédients et rend le flux d'ordres explicite. Le mid vaut $S_t=S_0+\sigma W_t$ (arithmétique, sur une session de trading). Le teneur de marché affiche un cours acheteur
$S^b_t=S_t-\delta^b_t$ et un cours vendeur $S^a_t=S_t+\delta^a_t$, où les profondeurs de cotation $\delta^b_t,\delta^a_t$ sont ses contrôles (symboles locaux de ce
chapitre~; $\delta$ seul désigne la fraction de calcul des intérêts). Les ordres de vente au marché atteignent le cours acheteur, et les achats au marché le cours vendeur, à des taux qui décroissent avec la profondeur.
Son inventaire $q_t$ augmente d'une unité à chaque exécution sur le cours acheteur et diminue d'une unité à chaque exécution sur le cours vendeur~; sa trésorerie $X_t$ varie des prix d'exécution.

\begin{definition}[Intensité d'exécution]\label{def:hf:inventory-models:intensity}
L'\emph{intensité d'exécution}\index{intensité d'exécution} (fill intensity) d'une cotation passive est le taux auquel elle est exécutée, en fonction de sa profondeur $\delta$ par rapport au
mid. Les modèles d'inventaire supposent $\lambda(\delta)=A\,e^{-k\delta}$, identique de part et d'autre, avec $A$ et $k$ constants (symboles locaux~: ici $A$ n'est pas une
annuité et $k$ n'est pas une log-moneyness ici).
\end{definition}

La forme exponentielle n'est pas une loi. Elle découle de l'hypothèse selon laquelle les ordres au marché arrivent à taux constant, avec des tailles dont l'impact sur le prix atteint une
profondeur $\delta$ avec la probabilité $e^{-k\delta}$~; elle est commode parce qu'elle rend le problème de contrôle soluble, et elle est testable, puisqu'un teneur de marché peut
mesurer à quelle fréquence elle est exécutée à chaque profondeur (voir la \cref{sec:hf:inventory-models:limits}).

\section{Le modèle d'Avellaneda--Stoikov}

\begin{definition}[Modèle d'Avellaneda--Stoikov]\label{def:hf:inventory-models:as}
Le \emph{modèle d'Avellaneda--Stoikov}\index{modèle d'Avellaneda--Stoikov} est le problème du teneur de marché avec un mid brownien arithmétique, des
\omterm{def:hf:inventory-models:intensity}{intensités d'exécution} exponentielles, un horizon $T$ et une utilité exponentielle de la richesse terminale~: le teneur de marché choisit $\delta^b_t,\delta^a_t$ pour maximiser
$\E\bigl[-\exp\bigl(-\gamma\,(X_T+q_TS_T)\bigr)\bigr]$, avec $\gamma>0$ son aversion au risque.
\end{definition}

\begin{sloppypar}L'équation de Hamilton--Jacobi--Bellman du livre~4, chapitre~9, s'applique. Avec l'ansatz $u(t,x,s,q)=-\exp(-\gamma(x+\theta(t,s,q)))$, les
conditions du premier ordre donnent les profondeurs optimales comme fonctions des différences de $\theta$ entre inventaires~; Avellaneda et Stoikov développent alors
$\theta$ à l'ordre deux en $q$ et obtiennent la solution approchée ci-dessous.\end{sloppypar}

\section{Prix de réservation et fourchette optimale}

\begin{definition}[Prix de réservation]\label{def:hf:inventory-models:reservation}
Le \emph{prix de réservation}\index{prix de réservation} d'un teneur de marché détenant $q$ est le prix auquel il est indifférent à sa position
courante et à une unité de plus ou de moins~; dans le \omterm{def:hf:inventory-models:as}{modèle d'Avellaneda--Stoikov}, approximativement
\[
S^{\mathrm{res}}_t=S_t-q_t\,\gamma\sigma^2\,(T-t).
\]
\end{definition}

\begin{proposition}[Les cotations d'Avellaneda--Stoikov]\label{prop:hf:inventory-models:as}
À l'ordre deux en inventaire, les cotations optimales sont centrées sur le \omterm{def:hf:inventory-models:reservation}{prix de réservation}, $S^{b,a}_t=S^{\mathrm{res}}_t\mp\tfrac12\Delta_t$, avec
une largeur totale
\[
\Delta_t=\delta^a_t+\delta^b_t=\gamma\sigma^2\,(T-t)+\frac{2}{\gamma}\ln\Bigl(1+\frac{\gamma}{k}\Bigr).
\]
\end{proposition}

\admitted
La démonstration figure dans l'article d'Avellaneda et Stoikov~; la \cref{prop:hf:inventory-models:cj} ci-dessous en dérive une solution exacte d'un modèle voisin.

Deux effets sont distingués. Le \omterm{def:hf:inventory-models:reservation}{prix de réservation} évolue à l'encontre de l'inventaire, proportionnellement à l'aversion au risque, à la variance et au temps restant~: un teneur de marché en position
longue affiche des cotations plus basses des deux côtés, vendant plus volontiers et achetant moins. La largeur comporte une composante de risque d'inventaire, $\gamma\sigma^2(T-t)$, et une
composante de monopole, $(2/\gamma)\ln(1+\gamma/k)$, qui tend vers $2/k$ lorsque $\gamma\to0$~: sans aversion au risque, le teneur de marché affiche encore une fourchette,
fixée par la vitesse à laquelle les exécutions décroissent avec la profondeur. Le skew des cotations est le décalage des cotations (skewing) présenté au livre~2, chapitre~15, mais dérivé ici d'un modèle plutôt que
d'une règle empirique.

\begin{definition}[Pénalité d'inventaire]\label{def:hf:inventory-models:penalty}
Une \emph{pénalité d'inventaire}\index{pénalité d'inventaire} remplace l'aversion au risque dans l'utilité par un coût prélevé sur la position tant qu'elle est détenue~: le
teneur de marché maximise $\E\bigl[X_T+q_TS_T-\alpha q_T^2-\phi\int_0^Tq_t^2\,dt\bigr]$, avec une pénalité continue $\phi$ et une pénalité terminale $\alpha$.
\end{definition}

La pénalité continue est une charge de risque que le desk choisit, et non une préférence qu'il a. Avec $\phi=\tfrac12\gamma\sigma^2$, elle pénalise exactement la variance
que la position ajoute par unité de temps~; un desk peut aussi la fixer à partir d'un budget de valeur en risque (VaR).

\begin{proposition}[Solution exacte avec une pénalité d'inventaire]\label{prop:hf:inventory-models:cj}
Avec un inventaire confiné à $\{-\bar q,\dots,\bar q\}$, on écrit la fonction valeur sous la forme $x+qs+h(t,q)$ et on pose $h=\tfrac1k\ln\omega$. Alors
$\omega(t)=\exp\bigl(M(T-t)\bigr)z$, avec $z_q=e^{-\alpha kq^2}$ et $M$ tridiagonale~: $M_{q,q}=-\phi kq^2$ et $Ae^{-1}$ à côté de la diagonale.
Les profondeurs optimales sont
\[
\delta^a(t,q)=\frac1k+h(t,q)-h(t,q-1),\qquad \delta^b(t,q)=\frac1k+h(t,q)-h(t,q+1),
\]
et le côté qui franchirait la borne d'inventaire n'est pas affiché.
\end{proposition}

\begin{proof}
L'équation de Hamilton--Jacobi--Bellman se réduit à $\partial_th-\phi q^2+A\sup_{\delta}e^{-k\delta}\bigl(\delta+h(t,q-1)-h(t,q)\bigr)+A\sup_\delta
e^{-k\delta}\bigl(\delta+h(t,q+1)-h(t,q)\bigr)=0$, avec $h(T,q)=-\alpha q^2$~; la diffusion du mid disparaît parce que la valeur est linéaire en $s$. La
condition du premier ordre du premier supremum donne $\delta^a=1/k+h(t,q)-h(t,q-1)$ et sa valeur $\tfrac1ke^{-1}e^{k(h(t,q-1)-h(t,q))}$, et de même
sur le cours acheteur. En substituant $h=\tfrac1k\ln\omega$, l'équation se ramène à $\partial_t\omega_q=\phi kq^2\omega_q-Ae^{-1}(\omega_{q-1}+\omega_{q+1})$, soit le
système linéaire $\partial_t\omega+M\omega=0$ avec $\omega(T)=z$.
\end{proof}

La \Cref{fig:hf:inventory-models:depths} montre la solution pour les paramètres du chapitre. À inventaire nul, le teneur de marché affiche des profondeurs symétriques au-dessus de $1/k$~;
en position longue, il abaisse la profondeur du cours vendeur et relève celle du cours acheteur, et pour de grands inventaires la profondeur du cours vendeur devient négative~: le modèle indique au teneur de marché de
coter sous le mid, ce qui, dans un carnet d'ordres réel, revient à croiser la fourchette. Le modèle ne connaît pas de fourchette à croiser~; le chapitre~4 borne plutôt
l'inventaire.

\begin{omfigure}
\begin{tikzpicture}
\begin{axis}[width=11cm,height=5.4cm,xlabel={\small inventaire $q$},ylabel={\small profondeur par rapport au mid},tick label style={font=\footnotesize},
  xmin=-10,xmax=10,ymin=-3,ymax=4.5,grid=major,grid style={gray!25},table/col sep=comma,
  legend style={font=\scriptsize,at={(0.5,-0.3)},anchor=north},legend columns=3]
\addplot[omDef,thick] table[x=q,y=b05]{figdata/hft/03-inventory-models/depths.csv};
\addplot[omDef,thick,dashed] table[x=q,y=a05]{figdata/hft/03-inventory-models/depths.csv};
\addplot[omThm,thick] table[x=q,y=b5]{figdata/hft/03-inventory-models/depths.csv};
\addplot[omThm,thick,dashed] table[x=q,y=a5]{figdata/hft/03-inventory-models/depths.csv};
\addplot[black,thick] table[x=q,y=b50]{figdata/hft/03-inventory-models/depths.csv};
\addplot[black,thick,dashed] table[x=q,y=a50]{figdata/hft/03-inventory-models/depths.csv};
\draw[gray] (axis cs:-10,0) -- (axis cs:10,0);
\legend{{achat, $\phi=0.5$},{vente, $\phi=0.5$},{achat, $\phi=5$},{vente, $\phi=5$},{achat, $\phi=50$},{vente, $\phi=50$}}
\end{axis}
\end{tikzpicture}

\omcaption{Profondeur optimale du cours acheteur (trait plein) et du cours vendeur (tirets) en début de séance, en fonction de l'inventaire, pour trois pénalités courantes $\phi$
(\cref{prop:hf:inventory-models:cj}~; $A=140$, $k=1.5$, $T=1$, $\alpha=1$, inventaire compris entre $\pm30$). Sous zéro, le cours vendeur est affiché sous le mid.
Données~: \texttt{firm.invmm.CJSolution}.}\label{fig:hf:inventory-models:depths}
\end{omfigure}

\section{L'expérience d'Avellaneda--Stoikov, reproduite}

Avellaneda et Stoikov ont simulé leurs cotations avec $S_0=100$, $T=1$, $\sigma=2$, des pas de temps de 0.005, $k=1.5$ et $A=140$, le mid se déplaçant vers le haut ou vers le bas
de $\sigma\sqrt{0.005}$ à chaque pas, et les ont comparées à une stratégie symétrique qui cote la même largeur moyenne autour du mid. La même
expérience, menée avec \texttt{firm.invmm} sur 1\,000 trajectoires avec les deux stratégies sur des nombres aléatoires communs, donne~:

\begin{center}
\small
\setlength{\tabcolsep}{4pt}%
\begin{tabular}{@{}>{\raggedright\arraybackslash}p{0.8cm}>{\raggedright\arraybackslash}p{1.9cm}>{\centering\arraybackslash}p{1.75cm}>{\centering\arraybackslash}p{1.75cm}>{\centering\arraybackslash}p{1.75cm}>{\centering\arraybackslash}p{1.75cm}>{\centering\arraybackslash}p{1.75cm}>{\centering\arraybackslash}p{1.75cm}@{}}
\toprule
$\gamma$ & stratégie & largeur & profit & écart-type (profit) & $q$ final & écart-type ($q$ final) & article~: profit, écart-type\\
\midrule
0.01 & inventaire & 1.35 & 68.45 & 9.06 & $-0.05$ & 5.24 & 68.6, 8.7\\
 & symétrique & 1.35 & 68.50 & 13.91 & $-0.05$ & 8.99 & 68.8, 12.8\\
0.1 & inventaire & 1.49 & 64.83 & 6.49 & 0.09 & 2.97 & 65.0, 6.6\\
 & symétrique & 1.49 & 68.15 & 14.00 & $-0.21$ & 8.66 & 68.4, 12.7\\
1 & inventaire & 3.03 & 31.44 & 4.83 & 0.05 & 1.66 & 31.4, 5.0\\
 & symétrique & 3.03 & 43.30 & 10.34 & $-0.14$ & 5.39 & 44.0, 11.0\\
\bottomrule
\end{tabular}
\end{center}

La reproduction concorde avec les tableaux publiés à l'erreur d'échantillonnage près. Au voisinage de la neutralité au risque, les deux stratégies rapportent la même chose et la
stratégie à inventaire est déjà moins variable~; pour $\gamma=0.1$, elle abandonne 5~\% du profit pour la moitié de l'écart-type~; pour $\gamma=1$, elle abandonne
27~\%. La \cref{fig:hf:inventory-models:path} retrace une trajectoire~: les cotations suivent le \omterm{def:hf:inventory-models:reservation}{prix de réservation}, qui s'écarte du mid chaque fois que l'inventaire
s'accumule.

\begin{omfigure}
\begin{tikzpicture}
\begin{axis}[width=11cm,height=5.4cm,xlabel={\small temps},ylabel={\small prix},tick label style={font=\footnotesize},xmin=0,xmax=1,
  grid=major,grid style={gray!25},table/col sep=comma,legend style={font=\scriptsize,at={(0.5,-0.3)},anchor=north},legend columns=4]
\addplot[black,thick] table[x=t,y=s]{figdata/hft/03-inventory-models/path.csv};
\addplot[omProp,thick,dashed] table[x=t,y=r]{figdata/hft/03-inventory-models/path.csv};
\addplot[omDef,thin] table[x=t,y=bid]{figdata/hft/03-inventory-models/path.csv};
\addplot[omThm,thin] table[x=t,y=ask]{figdata/hft/03-inventory-models/path.csv};
\legend{mid,prix de réservation,cours acheteur,cours vendeur}
\end{axis}
\end{tikzpicture}

\omcaption{Une trajectoire de l'expérience d'Avellaneda--Stoikov ($\gamma=0.1$)~: le mid, le \omterm{def:hf:inventory-models:reservation}{prix de réservation} et les cotations de la stratégie à inventaire. La
largeur se resserre à mesure que l'horizon approche. Données~: \texttt{hf\_inventory.one\_path}.}\label{fig:hf:inventory-models:path}
\end{omfigure}

Avec une pénalité courante au lieu de l'utilité, les cotations exactes de la \cref{prop:hf:inventory-models:cj} tracent une frontière de la moyenne en fonction de l'écart-type
du profit qui se situe nettement au-dessus de celle des cotations symétriques (\cref{fig:hf:inventory-models:frontier})~: une petite pénalité, $\phi=0.5$, conserve
97~\% du profit de la meilleure largeur symétrique (66.7 contre 68.5) avec la moitié de son écart-type (6.9 contre 13.9).

\begin{omfigure}
\begin{tikzpicture}
\begin{axis}[width=11cm,height=5.4cm,xlabel={\small écart-type du profit},ylabel={\small profit moyen},tick label style={font=\footnotesize},
  xmin=4,xmax=15,ymin=25,ymax=72,grid=major,grid style={gray!25},table/col sep=comma,
  legend style={font=\scriptsize,at={(0.5,-0.3)},anchor=north},legend columns=3]
\addplot[omThm,thick,mark=*,mark size=1.3pt] table[x=sd,y=profit]{figdata/hft/03-inventory-models/frontier_cj.csv};
\addplot[omDef,thick,dashed,mark=square*,mark size=1.3pt] table[x=sd,y=profit]{figdata/hft/03-inventory-models/frontier_sym.csv};
\addplot[black,only marks,mark=triangle*,mark size=2pt] table[x=sd,y=profit]{figdata/hft/03-inventory-models/frontier_as.csv};
\legend{{pénalité $\phi$ (exacte)},{largeur symétrique},{Avellaneda--Stoikov $\gamma$}}
\end{axis}
\end{tikzpicture}

\omcaption{Moyenne en fonction de l'écart-type du profit sur 1\,000 trajectoires de l'expérience d'Avellaneda--Stoikov~: les cotations exactes avec des pénalités courantes
$\phi$ de 0 à 50, des cotations symétriques de demi-largeurs de 0.5 à 2, et les trois aversions au risque du tableau. Données~:
\texttt{hf\_inventory.frontier}.}\label{fig:hf:inventory-models:frontier}
\end{omfigure}

\section{Les limites des formes closes}\label{sec:hf:inventory-models:limits}

Quatre hypothèses portent le modèle, et chacune échoue quelque part.

\begin{remark}[Ce que le modèle laisse de côté]\label{rem:hf:inventory-models:limits}
(i) \emph{Prix continus.} Les cotations vivent sur une grille de tick~; les profondeurs du modèle sont des nombres réels. (ii) \emph{Intensités exponentielles.} Les exécutions dépendent
de la position dans la file et du carnet d'ordres derrière la cotation, et non de la seule profondeur. (iii) \emph{Absence d'information.} Le mid est une martingale non affectée par les exécutions~:
il n'y a pas de sélection adverse, de sorte qu'une exécution n'est jamais une mauvaise nouvelle (le chapitre~4 l'ajoute). (iv) \emph{Paramètres constants.} La volatilité et le flux d'ordres varient
au cours de la journée et avec l'actualité.
\end{remark}

Les deux premières défaillances sont graves sur un marché à grand tick. Sur l'action simulée du livre~7, un teneur de marché explorateur, qui déplaçait ses cotations d'un lot au hasard
entre le meilleur cours et un, deux ou trois ticks derrière lui, toutes les cinq secondes pendant une heure, a passé environ 1\,640 secondes au meilleur cours et a obtenu 8 exécutions, et
environ 5\,280 secondes derrière lui sans en obtenir aucune~: la profondeur du carnet derrière le meilleur cours n'est jamais atteinte par les ordres au marché de ce marché. Le seul
choix consiste à se placer au meilleur cours ou à ne pas s'y placer.

Le problème reste un problème de contrôle. Avec $\lambda$ l'\omterm{def:hf:inventory-models:intensity}{intensité d'exécution} d'un lot au meilleur cours par côté et $c$ la capture par exécution, nette de
sélection adverse, la valeur $h(t,q)$ résout
\[
\partial_th-\phi q^2+\lambda\max\bigl(0,c+h(t,q-1)-h(t,q)\bigr)+\lambda\max\bigl(0,c+h(t,q+1)-h(t,q)\bigr)=0,
\]
et le teneur de marché se place d'un côté exactement lorsque la parenthèse de ce côté est positive (\cref{lst:hf:inventory-models:touch}). La cotation ne s'éloigne jamais du meilleur cours~;
le décalage des cotations du modèle continu devient une décision sur les côtés à afficher.

\section{Tutoriel~: du modèle au teneur de marché}\label{tut:hf:inventory-models:tut}

\begin{tutorial}
\textbf{Objectif.} Résoudre le modèle, reproduire l'expérience publiée, puis faire tourner la version limitée au meilleur cours sur le marché simulé. \textbf{État final~:} le
tableau ci-dessus, les \cref{fig:hf:inventory-models:depths,fig:hf:inventory-models:path,fig:hf:inventory-models:frontier}, et le tableau ci-dessous.
\begin{tutsteps}
\item \textbf{Résoudre.} \texttt{CJSolution} construit $M$, remonte $\omega$ depuis $T$ avec une seule exponentielle de matrice, en renormalisant à chaque pas (un
facteur constant disparaît de toutes les profondeurs, et sans lui $e^{-\alpha kq^2}$ tombe sous la précision machine à $q=30$).
\omcode{code/firm/invmm/firm_invmm.py}{55}{75}{La solution exacte avec une pénalité courante (\cref{prop:hf:inventory-models:cj}).}\label{lst:hf:inventory-models:cj}
\item \textbf{Reproduire.} \texttt{hf\_inventory.as\_table()} fait tourner les stratégies à inventaire et symétrique sur 1\,000 trajectoires communes pour chaque $\gamma$.
\item \textbf{Calibrer pour la bande.} Sur trois sessions de calibration de vingt minutes, le Quoter symétrique du chapitre~1, toujours au meilleur cours, a été
exécuté 241 fois~: $\lambda=0.033$ lot par seconde et par côté, et ses exécutions, valorisées dix secondes plus tard, donnent $c=0.139$ tick par action.
\item \textbf{Résoudre le problème limité au meilleur cours} à rebours sur une grille d'une seconde, puis le coter (\cref{lst:hf:inventory-models:touch}).
\omcode{code/firm/invmm/firm_invmm.py}{89}{110}{Le teneur de marché à grand tick, limité au meilleur cours~: il se place d'un côté seulement si sa parenthèse est positive.}\label{lst:hf:inventory-models:touch}
\item \textbf{Comparer} sur six sessions de vingt minutes du marché simulé avec \texttt{tape\_compare(phi)}.
\end{tutsteps}
\textbf{Pour aller plus loin.} Faire dépendre $c$ du déséquilibre du carnet d'ordres (chapitre~7)~; estimer $\lambda$ par moment de la journée.
\end{tutorial}

\begin{center}
\small
\begin{tabular}{@{}lccccc@{}}
\toprule
quoter (un lot, six sessions) & P\&L (\$) & sd (\$) & actions & $|q|$ moyen (lots) & $|q|$ max moyen\\
\midrule
symétrique, au plus 5 lots & $-65.17$ & 97.87 & 15\,783 & 2.51 & 5.00\\
limité au meilleur cours, $\phi=3\times10^{-4}$ & $-0.75$ & 29.53 & 14\,317 & 1.12 & 3.67\\
limité au meilleur cours, $\phi=10^{-3}$ & 11.00 & 8.41 & 11\,433 & 0.57 & 2.33\\
\bottomrule
\end{tabular}
\end{center}

Les quoters pénalisés échangent presque autant et détiennent une fraction de l'inventaire~; ils perdent moins et varient beaucoup moins. Sur six sessions, les moyennes sont
imprécises (l'erreur type du quoter symétrique est de 40~\$), et c'est la direction qui constitue la leçon plutôt que l'ampleur~: sur ce marché, un cours exécuté est
suivi en moyenne d'un mouvement adverse, si bien que l'inventaire accumulé par le quoter symétrique est précisément ce qui le fait perdre.

\section{Construction~: la bibliothèque de modèles d'inventaire}\label{bld:hf:inventory-models:invmm}

\begin{build}
\bfield{Rôle} Cotations sensibles à l'inventaire, issues de formes closes et de solutions exactes, et simulateur rapide pour comparer des politiques de cotation sur des
nombres aléatoires communs.
\bfield{Interface} \texttt{as\_quotes(s, q, tau, gamma, sigma, k)}, \texttt{as\_policy}~; \texttt{CJSolution(A, k, phi, alpha, qmax, T, n).depths(t, q)},
\texttt{cj\_policy}~; \texttt{TouchSolution(lam, c, phi, alpha, qmax, T, n).post(t, q)}~; \texttt{estimate\_intensity(depth, time, fills)}~;
\texttt{simulate(policy, T, dt, sigma, A, k, paths, seed, qmax, mid)}~; \texttt{symmetric(half)}.
\bfield{Règles} Les politiques ne voient que le temps et l'inventaire. Les nombres aléatoires du simulateur ne dépendent que de la graine et du pas, de sorte que les politiques
sont comparées sur les mêmes trajectoires. Un côté situé à la borne d'inventaire n'est jamais coté.
\bfield{Tests d'acceptation} \texttt{code/firm/invmm/tests/}~: la forme close à la main~; les profondeurs symétriques lorsque le stock est nul et les skews en miroir~; la profondeur de monopole $1/k$
loin de l'horizon sans pénalités~; l'\omterm{def:hf:inventory-models:intensity}{intensité d'exécution} retrouvée à partir de comptages synthétiques~; le simulateur déterministe et son nombre d'exécutions
égal à $2Ae^{-k\delta}T$~; la solution limitée au meilleur cours qui se place moins loin lorsque l'inventaire est élevé et ne dépasse jamais la borne.
\bfield{Pour aller plus loin} Une grille bidimensionnelle pour deux instruments corrélés (chapitre~4)~; des intensités dépendant de la position dans la file (chapitre~5).
\end{build}

\begin{omsources}
\item M.~Avellaneda et S.~Stoikov, ``High-frequency trading in a limit order book'', \emph{Quantitative Finance} 8(3), 2008.
\item T.~Ho et H.~R.~Stoll, ``Optimal dealer pricing under transactions and return uncertainty'', \emph{Journal of Financial Economics} 9(1), 1981.
\item Á.~Cartea, S.~Jaimungal et J.~Penalva, \emph{Algorithmic and High-Frequency Trading}, Cambridge University Press, 2015.
\end{omsources}

\section{Exercices}

\begin{exercise}[$\star$]\label{exo:hf:inventory-models:1}
Un teneur de marché avec $\gamma=0.1$, $\sigma=2$ et la moitié de la session restante ($T-t=0.5$) détient 3 unités à un mid de 100. Calculer son \omterm{def:hf:inventory-models:reservation}{prix de réservation}.
\end{exercise}

\begin{exercise}[$\star$]\label{exo:hf:inventory-models:2}
Avec $k=1.5$, calculer la largeur d'Avellaneda--Stoikov au début de la session ($T-t=1$) pour $\gamma=0.1$, ainsi que sa part de monopole.
\end{exercise}

\begin{exercise}[$\star$]\label{exo:hf:inventory-models:3}
Un cours affiché à une profondeur de 0.5 reste en place pendant 1\,000 secondes et est exécuté 40 fois~; à une profondeur de 1.0, il reste en place pendant 1\,000 secondes et est exécuté 18 fois. Estimer $A$ et $k$.
\end{exercise}

\begin{exercise}[$\star\star$]\label{exo:hf:inventory-models:4}
Montrer que la largeur d'Avellaneda--Stoikov tend vers $2/k$ lorsque $\gamma\to0$ à l'horizon, et expliquer pourquoi un teneur de marché neutre au risque cote
tout de même une fourchette.
\end{exercise}

\begin{exercise}[$\star\star$]\label{exo:hf:inventory-models:5}
À partir de la \cref{fig:hf:inventory-models:depths}, à quel inventaire la profondeur du cours vendeur devient-elle négative pour $\phi=5$, et que signifierait la cotation dans un
carnet d'ordres réel~?
\end{exercise}

\begin{exercise}[$\star\star$]\label{exo:hf:inventory-models:6}
Pourquoi la stratégie symétrique gagne-t-elle en moyenne plus que la stratégie à inventaire pour $\gamma=0.1$, et pourquoi cela n'est-il pas une raison de la préférer~?
\end{exercise}

\begin{exercise}[$\star\star\star$]\label{exo:hf:inventory-models:7}
\emph{Programmation.} Avec \texttt{CJSolution(140, 1.5, 5.0, 1.0, 30, 1.0, 200)}, donner les profondeurs au début de la session pour $q=0$ et pour $q=3$.
\end{exercise}

\begin{exercise}[$\star\star\star$]\label{exo:hf:inventory-models:8}
\emph{Trouver la faille.} «~Nous avons ajusté $A$ et $k$ sur nos cotations situées à un et deux ticks derrière le meilleur cours, sur un titre à un tick, et obtenu $k=0$~: les exécutions
ne dépendent pas de la profondeur, donc nous allons coter plus profond et gagner davantage.~»
\end{exercise}

\section{Problème~: trois achats d'affilée}

\begin{problem}[{Problème du week-end --- trois achats d'affilée}]\label{pb:hf:inventory-models:1}
Un teneur de marché vient d'acheter trois fois. Décider de ce qu'il faut coter ensuite, d'abord dans le modèle du manuel, puis dans le marché où il échange réellement.

\medskip\noindent\textbf{Partie I --- Le modèle.}
\begin{enumerate}
\item Énoncer le \omterm{def:hf:inventory-models:as}{modèle d'Avellaneda--Stoikov}~: dynamique, contrôles et objectif.
\item Définir l'\omterm{def:hf:inventory-models:intensity}{intensité d'exécution} et donner l'hypothèse du modèle.
\item Définir le \omterm{def:hf:inventory-models:reservation}{prix de réservation} et écrire sa forme approchée.
\item Écrire la largeur approchée et la décomposer en ses deux parties.
\item Calculer le \omterm{def:hf:inventory-models:reservation}{prix de réservation} et les cotations pour $q=3$, $\gamma=0.1$, $\sigma=2$, $k=1.5$, $T-t=0.5$, $S=100$.
\end{enumerate}

\medskip\noindent\textbf{Partie II --- La solution exacte.}
\begin{enumerate}[resume]
\item Définir une \omterm{def:hf:inventory-models:penalty}{pénalité d'inventaire} et relier la pénalité courante à l'aversion au risque.
\item Déduire les profondeurs de la \cref{prop:hf:inventory-models:cj} de l'équation de Hamilton--Jacobi--Bellman.
\item Pourquoi la volatilité du mid disparaît-elle de l'équation pour $h$~?
\item Donner les profondeurs pour $q=0$ et $q=3$ avec $\phi=5$, et interpréter la profondeur négative du cours vendeur pour un inventaire élevé.
\end{enumerate}

\medskip\noindent\textbf{Partie III --- L'expérience.}
\begin{enumerate}[resume]
\item Donner les paramètres de l'expérience.
\item Donner le profit et l'écart-type des deux stratégies pour les trois aversions au risque, et comparer avec l'article.
\item À quelle aversion au risque les profits des deux stratégies divergent-ils, et de combien~?
\item Qu'apporte la frontière des pénalités~?
\end{enumerate}

\medskip\noindent\textbf{Partie IV --- Le marché qu'il échange.}
\begin{enumerate}[resume]
\item Qu'a découvert le quoter explorateur derrière le meilleur cours, et pourquoi~?
\item Écrire l'équation limitée au meilleur cours et sa règle de décision.
\item Comment $\lambda$ et $c$ ont-ils été calibrés, et pourquoi sur des sessions distinctes~?
\item Énoncer le \emph{résultat nommé}~: le P\&L et l'inventaire des quoters symétrique et pénalisés sur le marché simulé, et, dans le modèle, le
profit et l'écart-type des deux stratégies pour $\gamma=0.1$.
\item Pourquoi six sessions suffisent-elles pour la direction et non pour l'ampleur~?
\item Quelle hypothèse du modèle le marché simulé viole-t-il le plus~?
\item En une phrase~: que doit faire un teneur de marché après trois achats d'affilée~?
\end{enumerate}
\end{problem}

\section{Questions d'entretien}

\begin{interviewq}[$\star$ \iqroles{trader}]\label{iq:hf:inventory-models:1}
Est-ce votre cours acheteur, votre cours vendeur ou les deux que vous déplacez, et dans quel sens~? Vous êtes long et voulez revenir à plat.
\end{interviewq}

\begin{interviewq}[$\star\star$ \iqroles{chercheur}]\label{iq:hf:inventory-models:2}
Dérivez la profondeur optimale d'un teneur de marché neutre au risque faisant face à une \omterm{def:hf:inventory-models:intensity}{intensité d'exécution} $Ae^{-k\delta}$, en faisant abstraction de l'inventaire.
\end{interviewq}

\begin{interviewq}[$\star\star$ \iqroles{chercheur}]\label{iq:hf:inventory-models:3}
Pourquoi le \omterm{def:hf:inventory-models:reservation}{prix de réservation} d'Avellaneda--Stoikov varie-t-il avec $T-t$, et est-ce raisonnable pour un teneur de marché qui échange chaque jour~?
\end{interviewq}

\begin{interviewq}[$\star\star$ \iqroles{développeur}]\label{iq:hf:inventory-models:4}
Votre solveur à exponentielle de matrice renvoie des profondeurs NaN pour de grands inventaires. Que s'est-il passé, et comment le corrigeriez-vous sans modifier la réponse~?
\end{interviewq}

\begin{interviewq}[$\star\star$ \iqroles{risque}]\label{iq:hf:inventory-models:5}
Comment choisiriez-vous la \omterm{def:hf:inventory-models:penalty}{pénalité d'inventaire} courante pour un desk doté d'une limite de perte quotidienne~?
\end{interviewq}

\begin{interviewq}[$\star\star\star$ \iqroles{chercheur}]\label{iq:hf:inventory-models:6}
Sur un titre à un tick, qu'est-ce qui remplace le choix de la profondeur, et comment formuleriez-vous le problème de contrôle~?
\end{interviewq}
